\documentclass[12pt,a4paper]{article}

\usepackage[a4paper,margin=18mm]{geometry}
\usepackage[T2A]{fontenc}
\usepackage[utf8]{inputenc}
\usepackage[russian]{babel}
\usepackage{amsmath,amssymb,mathtools}
\usepackage{array,booktabs,longtable}
\usepackage{xcolor}
\usepackage{hyperref}
\usepackage{tikz}

\hypersetup{
  colorlinks=true,
  linkcolor=blue!55!black,
  urlcolor=blue!55!black
}

\setlength{\parindent}{0pt}
\setlength{\parskip}{0.55em}
\setlength{\emergencystretch}{2em}
\renewcommand{\arraystretch}{1.25}
\newcolumntype{L}[1]{>{\raggedright\arraybackslash}p{#1}}

\begin{document}

% -------------------- Титульный лист --------------------
\begin{titlepage}
\centering
\vspace*{0.22\textheight}

{\LARGE\bfseries Ревью домашней работы\par}
\vspace{2.2cm}

{\large Автор: Лёвин Артём Александрович\par}
\vspace{0.8cm}

{\large 15 сентября 2026 года\par}

\vfill
\begin{tikzpicture}
  \draw[blue!35!black, line width=0.8pt]
    (0,0) .. controls (2.0,0.45) and (4.0,-0.45) .. (6.0,0)
          .. controls (8.0,0.45) and (10.0,-0.45) .. (12.0,0);
\end{tikzpicture}
\vspace{0.8cm}
\end{titlepage}

% -------------------- Сводная таблица --------------------
\newpage
\section*{Сводная таблица ошибок}

{\small
\setlength{\tabcolsep}{3pt}
\begin{longtable}{@{}L{1.4cm}L{3.0cm}L{5.0cm}L{2.2cm}L{2.6cm}@{}}
\toprule
\textbf{№ задания} & \textbf{Тема} & \textbf{Суть ошибки} & \textbf{Ваш ответ} & \textbf{Правильный ответ} \\
\midrule
\endfirsthead

\toprule
\textbf{№ задания} & \textbf{Тема} & \textbf{Суть ошибки} & \textbf{Ваш ответ} & \textbf{Правильный ответ} \\
\midrule
\endhead

\hyperref[task:16]{16} &
Вычитание десятичных дробей &
Решение задания №16 на приложенной работе не представлено. &
не указан &
$0{,}764$ \\[0.4em]

\hyperref[task:17]{17} &
Действия с обыкновенными дробями &
Решение задания №17 на приложенной работе не представлено. &
не указан &
$2$ \\[0.4em]

\hyperref[task:18]{18} &
Вычитание смешанных чисел &
Решение задания №18 на приложенной работе не представлено. &
не указан &
$1\frac{1}{14}$ \\[0.4em]

\hyperref[task:19]{19} &
Деление смешанного числа на произведение дробей &
Решение задания №19 на приложенной работе не представлено. &
не указан &
$5$ \\[0.4em]

\hyperref[task:20]{20} &
Сокращение дроби через простые множители &
Решение задания №20 на приложенной работе не представлено. &
не указан &
$\frac{5}{7}$ \\

\bottomrule
\end{longtable}
}

% -------------------- Разделитель --------------------
\newpage
\thispagestyle{empty}
\vspace*{\fill}
\begin{center}
{\LARGE\bfseries Разбор ошибок}
\end{center}
\vspace*{\fill}

% =========================================================
% Задание 16
% =========================================================
\newpage
\phantomsection
\label{task:16}
\section*{Задание 16}

\textbf{Текст задания.}
Вычислите:
\[
1{,}75-0{,}986.
\]

\textbf{Ваше решение.}
Решение задания №16 на приложенной работе не представлено.

\textbf{Ваш ответ.}
Не указан.

\textcolor{red}{\textbf{Ошибка:}} задание не выполнено.

\textbf{Где именно возникает ошибка.}
Последний корректный шаг отсутствует, потому что решение задания №16 не начато.
Первый неполный шаг --- отсутствие вычисления разности $1{,}75-0{,}986$ и итогового ответа.

\textbf{Почему это неверно.}
В задании требуется вычислить разность двух десятичных дробей.
Чтобы выполнить вычитание без ошибки, нужно записать числа так, чтобы запятые стояли друг под другом,
а одинаковые разряды --- единицы, десятые, сотые и тысячные --- находились в одних столбцах.
Число $1{,}75$ удобно дописать как $1{,}750$: добавление нуля справа после последнего десятичного знака значение числа не меняет.

\textcolor{blue!60!black}{\textbf{Ключевая идея:}}
перед вычитанием десятичных дробей выровнять запятые и при необходимости дописать справа нули, чтобы число десятичных разрядов совпало.

\textbf{Исправление.}
Запишем первое число с тремя знаками после запятой:
\[
1{,}75=1{,}750.
\]
Тогда вычисляем
\[
1{,}750-0{,}986.
\]

\textbf{Полное правильное решение.}
Выполним вычитание по разрядам:
\[
1{,}750-0{,}986=0{,}764.
\]
Можно проверить вычисление сложением:
\[
0{,}764+0{,}986=1{,}750=1{,}75.
\]

\textbf{Проверка.}
Проверочное сложение возвращает уменьшаемое, поэтому разность найдена верно.
Кроме того, результат должен быть положительным и меньше $1{,}75$, что также выполняется.

\textcolor{teal!60!black}{\textbf{Итог:}}
\[
\boxed{0{,}764}.
\]

\textbf{Задача на закрепление.}
Вычислите:
\[
2{,}43-1{,}587.
\]

% =========================================================
% Задание 17
% =========================================================
\newpage
\phantomsection
\label{task:17}
\section*{Задание 17}

\textbf{Текст задания.}
Вычислите:
\[
\left(\frac{3}{5}+\frac{7}{10}\right):\frac{13}{20}.
\]

\textbf{Ваше решение.}
Решение задания №17 на приложенной работе не представлено.

\textbf{Ваш ответ.}
Не указан.

\textcolor{red}{\textbf{Ошибка:}} задание не выполнено.

\textbf{Где именно возникает ошибка.}
Последний корректный шаг отсутствует: решение задания №17 не начато.
Первый неполный шаг --- не выполнено действие в скобках, после которого нужно разделить полученную дробь на $\frac{13}{20}$.

\textbf{Почему это неверно.}
В выражении сначала выполняется действие в скобках.
Дроби $\frac{3}{5}$ и $\frac{7}{10}$ имеют разные знаменатели, поэтому перед сложением их нужно привести к общему знаменателю.
После этого деление на дробь заменяется умножением на обратную дробь.

\textcolor{blue!60!black}{\textbf{Ключевая идея:}}
сначала выполнить скобки, затем заменить деление на дробь умножением на обратную дробь и сократить общие множители.

\textbf{Исправление.}
Приведём дроби в скобках к знаменателю $10$:
\[
\frac{3}{5}=\frac{6}{10}.
\]
Тогда
\[
\frac{6}{10}+\frac{7}{10}=\frac{13}{10}.
\]
Теперь нужно вычислить
\[
\frac{13}{10}:\frac{13}{20}.
\]

\textbf{Полное правильное решение.}
\[
\left(\frac{3}{5}+\frac{7}{10}\right):\frac{13}{20}
=\left(\frac{6}{10}+\frac{7}{10}\right):\frac{13}{20}
=\frac{13}{10}:\frac{13}{20}.
\]
Заменяем деление умножением на обратную дробь:
\[
\frac{13}{10}\cdot\frac{20}{13}.
\]
Сокращаем множитель $13$ и число $20$ с $10$:
\[
\frac{13}{10}\cdot\frac{20}{13}=2.
\]

\textbf{Проверка.}
Если результат равен $2$, то при умножении делителя на частное должно получиться значение скобок:
\[
\frac{13}{20}\cdot 2=\frac{13}{10}.
\]
Это совпадает с найденной суммой в скобках.

\textcolor{teal!60!black}{\textbf{Итог:}}
\[
\boxed{2}.
\]

\textbf{Задача на закрепление.}
Вычислите:
\[
\left(\frac{1}{2}+\frac{3}{4}\right):\frac{5}{8}.
\]

% =========================================================
% Задание 18
% =========================================================
\newpage
\phantomsection
\label{task:18}
\section*{Задание 18}

\textbf{Текст задания.}
Вычислите:
\[
2\frac{3}{7}-1\frac{5}{14}.
\]

\textbf{Ваше решение.}
Решение задания №18 на приложенной работе не представлено.

\textbf{Ваш ответ.}
Не указан.

\textcolor{red}{\textbf{Ошибка:}} задание не выполнено.

\textbf{Где именно возникает ошибка.}
Последний корректный шаг отсутствует, поскольку решение задания №18 не начато.
Первый неполный шаг --- не приведены дробные части смешанных чисел к общему знаменателю и не выполнено вычитание.

\textbf{Почему это неверно.}
У дробных частей знаменатели $7$ и $14$.
Удобно привести $\frac{3}{7}$ к знаменателю $14$:
\[
\frac{3}{7}=\frac{6}{14}.
\]
После этого можно вычесть отдельно целые и дробные части, потому что $\frac{6}{14}$ больше $\frac{5}{14}$ и занимать единицу не требуется.

\textcolor{blue!60!black}{\textbf{Ключевая идея:}}
при вычитании смешанных чисел сначала привести дробные части к общему знаменателю, а затем проверить, можно ли вычесть дробную часть вычитаемого без займа единицы.

\textbf{Исправление.}
\[
2\frac{3}{7}=2\frac{6}{14}.
\]
Поэтому
\[
2\frac{6}{14}-1\frac{5}{14}
=(2-1)+\left(\frac{6}{14}-\frac{5}{14}\right).
\]

\textbf{Полное правильное решение.}
\[
2\frac{3}{7}-1\frac{5}{14}
=2\frac{6}{14}-1\frac{5}{14}
=1\frac{1}{14}.
\]

\textbf{Проверка.}
Прибавим вычитаемое к найденной разности:
\[
1\frac{1}{14}+1\frac{5}{14}
=2\frac{6}{14}
=2\frac{3}{7}.
\]
Получили уменьшаемое.

\textcolor{teal!60!black}{\textbf{Итог:}}
\[
\boxed{1\frac{1}{14}}.
\]

\textbf{Задача на закрепление.}
Вычислите:
\[
4\frac{5}{6}-2\frac{7}{12}.
\]

% =========================================================
% Задание 19
% =========================================================
\newpage
\phantomsection
\label{task:19}
\section*{Задание 19}

\textbf{Текст задания.}
Вычислите:
\[
3\frac{1}{3}:\left(\frac{5}{6}\cdot\frac{4}{5}\right).
\]

\textbf{Ваше решение.}
Решение задания №19 на приложенной работе не представлено.

\textbf{Ваш ответ.}
Не указан.

\textcolor{red}{\textbf{Ошибка:}} задание не выполнено.

\textbf{Где именно возникает ошибка.}
Последний корректный шаг отсутствует: решение задания №19 не начато.
Первый неполный шаг --- не вычислено произведение в скобках и не выполнено последующее деление смешанного числа на полученную дробь.

\textbf{Почему это неверно.}
По порядку действий сначала нужно вычислить произведение в скобках.
Затем смешанное число $3\frac{1}{3}$ переводится в неправильную дробь.
После этого деление на дробь заменяется умножением на обратную дробь.

\textcolor{blue!60!black}{\textbf{Ключевая идея:}}
соблюдать порядок действий: сначала скобки; смешанное число перевести в неправильную дробь; деление заменить умножением на обратную дробь.

\textbf{Исправление.}
Вычислим скобки:
\[
\frac{5}{6}\cdot\frac{4}{5}=\frac{4}{6}=\frac{2}{3}.
\]
Переведём смешанное число:
\[
3\frac{1}{3}=\frac{10}{3}.
\]
Получаем
\[
\frac{10}{3}:\frac{2}{3}.
\]

\textbf{Полное правильное решение.}
\[
3\frac{1}{3}:\left(\frac{5}{6}\cdot\frac{4}{5}\right)
=\frac{10}{3}:\frac{2}{3}
=\frac{10}{3}\cdot\frac{3}{2}
=5.
\]

\textbf{Проверка.}
Умножим найденное частное на значение скобок:
\[
5\cdot\frac{2}{3}=\frac{10}{3}=3\frac{1}{3}.
\]
Получили исходное делимое.

\textcolor{teal!60!black}{\textbf{Итог:}}
\[
\boxed{5}.
\]

\textbf{Задача на закрепление.}
Вычислите:
\[
2\frac{1}{4}:\left(\frac{3}{5}\cdot\frac{5}{6}\right).
\]

% =========================================================
% Задание 20
% =========================================================
\newpage
\phantomsection
\label{task:20}
\section*{Задание 20}

\textbf{Текст задания.}
Сократите дробь
\[
\frac{360}{504},
\]
предварительно разложив числитель и знаменатель на простые множители.

\textbf{Ваше решение.}
Решение задания №20 на приложенной работе не представлено.

\textbf{Ваш ответ.}
Не указан.

\textcolor{red}{\textbf{Ошибка:}} задание не выполнено; обязательное разложение на простые множители отсутствует.

\textbf{Где именно возникает ошибка.}
Последний корректный шаг отсутствует, потому что решение задания №20 не начато.
Первый неполный шаг --- не выполнено требуемое разложение числителя $360$ и знаменателя $504$ на простые множители.

\textbf{Почему это неверно.}
В условии указано не только сократить дробь, но и предварительно разложить числитель и знаменатель на простые множители.
Значит, даже верный сокращённый ответ без такого разложения не выполнил бы условие полностью.
После разложения одинаковые простые множители в числителе и знаменателе можно сократить.

\textcolor{blue!60!black}{\textbf{Ключевая идея:}}
разложить оба числа на простые множители, выделить общую часть разложений и сократить только одинаковые множители числителя и знаменателя.

\textbf{Исправление.}
Разложим числитель:
\[
360=36\cdot10=2^2\cdot3^2\cdot2\cdot5=2^3\cdot3^2\cdot5.
\]
Разложим знаменатель:
\[
504=63\cdot8=3^2\cdot7\cdot2^3=2^3\cdot3^2\cdot7.
\]

\textbf{Полное правильное решение.}
\[
\frac{360}{504}
=\frac{2^3\cdot3^2\cdot5}{2^3\cdot3^2\cdot7}.
\]
Сокращаем общие множители $2^3$ и $3^2$:
\[
\frac{360}{504}=\frac{5}{7}.
\]

\textbf{Проверка.}
Общий сокращённый множитель равен
\[
2^3\cdot3^2=8\cdot9=72.
\]
Проверяем:
\[
360:72=5,
\qquad
504:72=7.
\]
Следовательно, исходная дробь действительно сокращается до $\frac{5}{7}$.

\textcolor{teal!60!black}{\textbf{Итог:}}
\[
\boxed{\frac{5}{7}}.
\]

\textbf{Задача на закрепление.}
Сократите дробь
\[
\frac{420}{588},
\]
предварительно разложив числитель и знаменатель на простые множители.

\vfill
\begin{center}
\small Лёвин Артём Александрович эксклюзивно для Ксении Васильченко
\end{center}

\end{document}
